\ifdefined\gkver\else\def\gkver{student}\fi
\documentclass[\gkver]{gaokaozhenti}
\newcommand{\ccnum}[1]{{\fontspec{Noto Sans CJK JP}#1}}
\begin{document}

\chapter{2018年全国理综Ⅱ}
%% paperType: 理综物理部分

\begin{enumerate}
\item
%% number: 14
%% paperName: 2018年全国理综Ⅱ第 14 题 6 分
%% typeId: 1
%% type: 单选题
%% chapter: 功和能
%% point: 动能定理
%% method: 无
%% score: 6
%% degree: 850
%% duplicateId: 0
%% body:
如图，某同学用绳子拉动木箱，使它从静止开始沿粗糙水平路面运动至具有某一速度，木箱获得的动能一定\xzanswer{A}

\onepicture[h=3.2cm]{figs/14.png}

\fourchoices[answer=A]
{小于拉力所做的功}
{等于拉力所做的功}
{等于克服摩擦力所做的功}
{大于克服摩擦力所做的功}

%% answer: A
%% memo:
\memoanswer{
木箱由静止开始沿粗糙水平路面运动至具有某一速度的过程中，拉力和摩擦力做功，对整个过程应用动能定理，得 $W_{\text{拉}}-W_{f}=\Delta E_{k}=\frac{1}{2}mv^{2}-0$，可知木箱获得的动能小于拉力所做的功，与克服摩擦力做功的大小关系不确定，故 A 正确。
}

\item
%% number: 15
%% paperName: 2018年全国理综Ⅱ第 15 题 6 分
%% typeId: 1
%% type: 单选题
%% chapter: 运动和力
%% point: 牛顿定律的应用
%% method: 近似与估算
%% score: 6
%% degree: 750
%% duplicateId: 0
%% body:
高空坠物极易对行人造成伤害。若一个 $50\Ug$ 的鸡蛋从一居民楼的 25 层坠下，与地面的撞击时间约为 $2\UmsT$，则该鸡蛋对地面产生的冲击力约为\xzanswer{C}

\fourchoices[answer=C]
{$10\UN$}
{$10^{2}\UN$}
{$10^{3}\UN$}
{$10^{4}\UN$}

%% answer: C
%% memo:
\memoanswer{
鸡蛋由 25 层楼坠下，由于居民楼每层高约 $3\Um$，故下落高度 $h=3\times24\Um=72\Um$，鸡蛋做自由落体运动，由动能定理可知 $mgh=\frac{1}{2}mv^{2}$，解得 $v\approx38\Ums$。鸡蛋与地面碰撞后速度变为零，取向下为正方向，由动量定理得 $(mg-F_{N})t=0-mv$，解得 $F_{N}\approx950\UN$，由牛顿第三定律可知鸡蛋对地面产生的冲击力约为 $10^{3}\UN$，故 C 正确。
}

\item
%% number: 16
%% paperName: 2018年全国理综Ⅱ第 16 题 6 分
%% typeId: 1
%% type: 单选题
%% chapter: 曲线运动
%% point: 天体运动
%% method: 无
%% score: 6
%% degree: 550
%% duplicateId: 0
%% body:
2018 年 2 月，我国 $500\Um$ 口径射电望远镜（天眼）发现毫秒脉冲星“J0318+0253”，其自转周期 $T=5.19\UmsT$，假设星体为质量均匀分布的球体，已知万有引力常量为 $6.67\times10^{-11}\UN\cdot\Um/\Ukg^{2}$。以周期 $T$ 稳定自转的星体的密度最小值约为\xzanswer{C}

\fourchoices[answer=C]
{$5\times10^{9}\Ukgmc$}
{$5\times10^{12}\Ukgmc$}
{$5\times10^{15}\Ukgmc$}
{$5\times10^{18}\Ukgmc$}

%% answer: C
%% memo:
\memoanswer{
假设星体表面有一质量为 $m$ 的物体，可随星体稳定自转且对星体表面无压力，由万有引力定律可知 $G\frac{Mm}{R^{2}}=m\frac{4\pi^{2}}{T^{2}}R$，又因为 $M=\rho\cdot\frac{4}{3}\pi R^{3}$，可知最小密度 $\rho=\frac{3\pi}{GT^{2}}\approx5\times10^{15}\Ukgmc$，故 C 正确。
}

\item
%% number: 17
%% paperName: 2018年全国理综Ⅱ第 17 题 6 分
%% typeId: 1
%% type: 单选题
%% chapter: 光学
%% point: 光电效应
%% method: 无
%% score: 6
%% degree: 550
%% duplicateId: 0
%% body:
用波长为 $300\Unm$ 的光照射锌板，电子逸出锌板表面的最大初动能为 $1.28\times10^{-19}\UJ$。已知普朗克常量为 $6.63\times10^{-34}\UJcdots$，真空中的光速为 $3.00\times10^{8}\Ums$，能使锌产生光电效应的单色光的最低频率约为\xzanswer{B}

\fourchoices[answer=B]
{$1\times10^{14}\UHz$}
{$8\times10^{14}\UHz$}
{$2\times10^{15}\UHz$}
{$8\times10^{15}\UHz$}

%% answer: B
%% memo:
\memoanswer{
由光电效应方程得 $E_{\mathrm{km}}=h\nu-W_{0}$，即 $W_{0}=h\nu-E_{\mathrm{km}}$，而 $W_{0}=h\nu_{c}$，联立解得 $\nu_{c}=\nu-\frac{E_{\mathrm{km}}}{h}=\frac{c}{\lambda}-\frac{E_{\mathrm{km}}}{h}\approx8\times10^{14}\UHz$，故 B 正确。
}

\item
%% number: 18
%% paperName: 2018年全国理综Ⅱ第 18 题 6 分
%% typeId: 1
%% type: 单选题
%% chapter: 电磁感应
%% point: 电磁感应中图象
%% method: 无
%% score: 6
%% degree: 450
%% duplicateId: 0
%% body:
如图，在同一平面内有两根平行长导轨，导轨间存在依次相邻的矩形匀强磁场区域，区域宽度均为 $l$，磁感应强度大小相等、方向交替向上向下。一边长为 $\frac{3}{2}l$ 的正方形金属线框在导轨上向左匀速运动，线框中感应电流 $i$ 随时间 $t$ 变化的正确图线可能是\xzanswer{D}

\onepicture[h=3.5cm]{figs/18.png}

\fourchoices[answer=D, ispicture=true, h=2.2cm]
{figs/18a.png}
{figs/18b.png}
{figs/18c.png}
{figs/18d.png}

%% answer: D
%% memo:
\memoanswer{
线框由图示位置向左匀速移动 $\frac{l}{2}$ 的过程中，金属线框左、右两边切割磁感线，由右手定则可知其产生的感应电流沿顺时针方向，$i=\frac{2Blv}{R}$，电流大小恒定；在线框移动 $\frac{l}{2}\sim l$ 过程中，左、右两边产生的感应电动势大小相等，方向相反，感应电流为零；在线框移动 $l\sim\frac{3}{2}l$ 过程中，金属线框左、右两边切割磁感线，由右手定则可知其产生的感应电流沿逆时针方向，电流大小恒定；在线框移动 $\frac{3}{2}l\sim2l$ 过程中，左、右两边产生的感应电动势大小相等，方向相反，感应电流为零，之后周期性重复，故 D 正确。
}

\item
%% number: 19
%% paperName: 2018年全国理综Ⅱ第 19 题 6 分
%% typeId: 2
%% type: 多选题
%% chapter: 直线运动
%% point: 运动图象
%% method: 无
%% score: 6
%% degree: 650
%% duplicateId: 0
%% body:
甲、乙两汽车在同一条平直公路上同向运动，其速度-时间图像分别如图中甲、乙两条曲线所示。已知两车在 $t_{2}$ 时刻并排行驶，下列说法正确的是\xzanswer{BD}

\onepicture[h=3.6cm]{figs/19.png}

\fourchoices[answer=BD]
{两车在 $t_{1}$ 时刻也并排行驶}
{$t_{1}$ 时刻甲车在后，乙车在前}
{甲车的加速度大小先增大后减小}
{乙车的加速度大小先减小后增大}

%% answer: BD
%% memo:
\memoanswer{
已知 $t_{2}$ 时刻，甲、乙两车并排行驶，由 $v-t$ 图像与时间轴围成图形的面积表示位移可知，在 $t_{1}\sim t_{2}$ 时间内，甲的位移大于乙的位移，故在 $t_{1}$ 时刻，甲车在乙车的后方，故 A 错误、B 正确；$v-t$ 图像的斜率表示物体运动的加速度大小，由图像可知，甲、乙两车的加速度大小均先减小后增大，故 C 错误、D 正确。
}

\item
%% number: 20
%% paperName: 2018年全国理综Ⅱ第 20 题 6 分
%% typeId: 2
%% type: 多选题
%% chapter: 磁场
%% point: 磁场的叠加
%% method: 无
%% score: 6
%% degree: 450
%% duplicateId: 0
%% body:
如图，纸面内有两条互相垂直的长直绝缘导线 $L_{1}$、$L_{2}$，$L_{1}$ 中的电流方向向左，$L_{2}$ 中的电流方向向上；$L_{1}$ 的正上方有 a、b 两点，它们相对于 $L_{2}$ 对称。整个系统处于匀强外磁场中，外磁场的磁感应强度大小为 $B_{0}$，方向垂直于纸面向外。已知 a、b 两点的磁感应强度大小分别为 $\frac{1}{3}B_{0}$ 和 $\frac{1}{2}B_{0}$，方向也垂直于纸面向外。则\xzanswer{AC}

\onepicture[w=4.3cm]{\includesvg{figs/20}}

\fourchoices[answer=AC]
{流经 $L_{1}$ 的电流在 b 点产生的磁感应强度大小为 $\frac{7}{12}B_{0}$}
{流经 $L_{1}$ 的电流在 a 点产生的磁感应强度大小为 $\frac{1}{12}B_{0}$}
{流经 $L_{2}$ 的电流在 b 点产生的磁感应强度大小为 $\frac{1}{12}B_{0}$}
{流经 $L_{2}$ 的电流在 a 点产生的磁感应强度大小为 $\frac{7}{12}B_{0}$}

%% answer: AC
%% memo:
\memoanswer{
设导线 $L_{1}$ 在 a、b 两点产生的磁感应强度大小为 $B_{1}$，导线 $L_{2}$ 在 a、b 两点产生的磁感应强度大小为 $B_{2}$，以垂直于纸面向外为正方向，在 a 点有 $B_{0}-B_{1}-B_{2}=\frac{1}{3}B_{0}$，在 b 点有 $B_{0}-B_{1}+B_{2}=\frac{1}{2}B_{0}$，解得 $B_{1}=\frac{7}{12}B_{0}$，$B_{2}=\frac{1}{12}B_{0}$，故 AC 正确。
}

\item
%% number: 21
%% paperName: 2018年全国理综Ⅱ第 21 题 6 分
%% typeId: 2
%% type: 多选题
%% chapter: 电场
%% point: 电势能电势差电场力做功
%% method: 无
%% score: 6
%% degree: 450
%% duplicateId: 0
%% body:
如图，同一平面内的 a、b、c、d 四点处于匀强电场中，电场方向与此平面平行，M 为 a、c 连线的中点，N 为 b、d 连线的中点。一电荷量为 $q$（$q>0$）的粒子从 a 点移动到 b 点，其电势能减小 $W_{1}$；若该粒子从 c 点移动到 d 点，其电势能减小 $W_{2}$，下列说法正确的是\xzanswer{BD}

\onepicture[w=4.3cm]{\includesvg{figs/21}}

\fourchoices[answer=BD]
{此匀强电场的场强方向一定与 a、b 两点连线平行}
{若该粒子从 M 点移动到 N 点，则电场力做功一定为 $\frac{W_{1}+W_{2}}{2}$}
{若 c、d 之间的距离为 $L$，则该电场的场强大小一定为 $\frac{W_{2}}{qL}$}
{若 $W_{1}=W_{2}$，则 a、M 两点之间的电势差一定等于 b、N 两点之间的电势差}

%% answer: BD
%% memo:
\memoanswer{
分析可知，此匀强电场方向不一定与 a、b 连线平行，故 A 错误；由题可知，粒子从 a 移动到 b 过程中，$W_{1}=qU_{ab}$，由 c 移动到 d 过程中，$W_{2}=qU_{cd}$，设该粒子从 M 移动到 N 电场力做功为 $W$，$W=qU_{MN}$，$U_{MN}=\varphi_{M}-\varphi_{N}$，在匀强电场中，由于 M、N 分别为 a、c 连线与 b、d 连线的中点，故 $\varphi_{M}=\frac{\varphi_{a}+\varphi_{c}}{2}$，$\varphi_{N}=\frac{\varphi_{b}+\varphi_{d}}{2}$，$U_{MN}=\varphi_{M}-\varphi_{N}=\frac{U_{ab}+U_{cd}}{2}$，即 $W=\frac{W_{1}+W_{2}}{2}$，故 B 正确；由题可知，匀强电场方向大致由左向右，无法判断其方向是否与 c、d 连线平行，则该电场的场强大小不一定为 $\frac{W_{2}}{qL}$，故 C 错误；$U_{aM}=\varphi_{a}-\varphi_{M}=\varphi_{a}-\frac{\varphi_{a}+\varphi_{c}}{2}=\frac{\varphi_{a}-\varphi_{c}}{2}$，$U_{bN}=\varphi_{b}-\varphi_{N}=\varphi_{b}-\frac{\varphi_{b}+\varphi_{d}}{2}=\frac{\varphi_{b}-\varphi_{d}}{2}$，若 $W_{1}=W_{2}$，则 $U_{ab}=U_{cd}$，即 $\varphi_{a}-\varphi_{b}=\varphi_{c}-\varphi_{d}$，得 $\varphi_{a}-\varphi_{c}=\varphi_{b}-\varphi_{d}$，即 $U_{aM}=U_{bN}$，故 D 正确。
}

\item
%% number: 22
%% paperName: 2018年全国理综Ⅱ第 22 题 6 分
%% typeId: 4
%% type: 实验题
%% chapter: 电路
%% point: 电表改装
%% method: 无
%% score: 6
%% degree: 750
%% duplicateId: 0
%% body:
某同学组装一个多用电表。可用的器材有：微安表头（量程 $100\UmuA$，内阻 $900\UO$）；电阻箱 $R_{1}$（阻值范围 $0\sim999.9\UO$）；电阻箱 $R_{2}$（阻值范围 $0\sim99999.9\UO$）；导线若干。

要求利用所给器材先组装一个量程为 $1\UmA$ 的直流电流表，在此基础上再将它改装成量程为 $3\UV$ 的直流电压表。组装好的多用电表有电流 $1\UmA$ 和电压 $3\UV$ 两挡。

回答下列问题：
\begin{enumerate}
	\item
	在虚线框内画出电路图并标出 $R_{1}$ 和 $R_{2}$，其中 $*$ 为公共接线柱，a 和 b 分别是电流挡和电压挡的接线柱\tkanswer[2cm]{}。
	\item
	电阻箱的阻值 $R_{1}=$\tkanswer[1.6cm]{$100$} $\UO$；$R_{2}=$\tkanswer[1.8cm]{$2910$} $\UO$（保留到个位）。
\end{enumerate}

\onepicture[h=3.2cm, align=right]{figs/22.png}

%% answer:
\jdanswer{
\begin{enumerate}
	\item
	电路图见详解

	\item
	$100$，$2910$

\end{enumerate}
}
%% memo:
\memoanswer{
（1）微安表头改装为电流表，应与电阻箱并联分流，分流电阻较小，故与微安表头并联的电阻箱选 $R_{1}$，再将改装后的电流表与电阻箱串联分压，改装成电压表，分压电阻较大，则与电流表串联的电阻箱选 $R_{2}$，电路连接如图所示。

\onepicture[w=6.5cm]{figs/22_an.png}

（2）设微安表头满偏时分流电阻中电流为 $I_{1}$，则有 $I_{1}R_{1}=I_{g}R_{g}$，$I_{1}+I_{g}=1\UmA$，可得 $R_{1}=100\UO$。设微安表头满偏时 $R_{2}$ 两端电压为 $U_{2}$，则有 $U_{1}=I_{g}R_{g}$，$U_{1}+U_{2}=3\UV$，$U_{2}=(I_{1}+I_{g})R_{2}$，可得 $R_{2}=2910\UO$。
}

\item
%% number: 23
%% paperName: 2018年全国理综Ⅱ第 23 题 9 分
%% typeId: 4
%% type: 实验题
%% chapter: 力物体的平衡
%% point: 摩擦力
%% method: 无
%% score: 9
%% degree: 450
%% duplicateId: 0
%% body:
某同学用图\ref{2018全国理综Ⅱ23a}所示的装置测量木块与木板之间的摩擦因数。跨过光滑定滑轮的细线两端分别与木块和弹簧秤相连，滑轮和木块之间的细线保持水平，在木块上放置砝码。缓慢向左拉动水平放置的木板，当木块和砝码相对桌面静止且木板仍在继续滑动时，弹簧秤的示数即为木块受到的滑动摩擦力的大小。某次实验所得数据在下表中给出，其中 $f_{4}$ 的值从图\ref{2018全国理综Ⅱ23b}中弹簧秤的示数读出。

\begin{table}[htbp]
\centering
\begin{tblr}{|c|c|c|c|c|c|}
\hline
砝码的质量 $m/\Ukg$ & 0.05 & 0.10 & 0.15 & 0.20 & 0.25 \\
\hline
滑动摩擦力 $f/\UN$ & 2.15 & 2.36 & 2.55 & $f_{4}$ & 2.93 \\
\hline
\end{tblr}
\end{table}

回答下列问题：
\begin{enumerate}
	\item
	$f_{4}=$\tkanswer[1.6cm]{$2.75$} $\UN$；
	\item
	在图\ref{2018全国理综Ⅱ23c}的坐标纸上补齐未画出的数据点并绘出 $f-m$ 图线\tkanswer[2cm]{}；
	\item
	$f$ 与 $m$、木块质量 $M$、木板和木块之间的滑动摩擦因数 $\mu$ 及重力加速度大小 $g$ 之间的关系式 $f=$\tkanswer[3cm]{$\mu(M+m)g$}，$f-m$ 图线（直线）的斜率的表达式 $k=$\tkanswer[2cm]{$\mu g$}；
	\item
	取 $g=9.80\Umsq$，由绘出的 $f-m$ 图线求得 $\mu=$\tkanswer[1.6cm]{$0.40$}（保留 2 位有效数字）。
\end{enumerate}

\threepicture[numA=a, numB=b, numC=c, labelA=2018全国理综Ⅱ23a, labelB=2018全国理综Ⅱ23b, labelC=2018全国理综Ⅱ23c, wA=4.5cm, wB=2.5cm, wC=5.0cm, align=right]
{figs/23a.png}
{figs/23b.png}
{\includesvg{figs/23c}}

%% answer:
\jdanswer{
\begin{enumerate}
	\item
	$2.75$

	\item
	如图所示

	\item
	$\mu(M+m)g$，$\mu g$

	\item
	$0.40$

\end{enumerate}
}
%% memo:
\memoanswer{
（1）弹簧秤读数时应估读，弹簧秤示数为 $2.75\UN$。

（2）按数据描点连线如图所示。

\onepicture[w=7.0cm]{\includesvg{figs/23_an}}

（3）由于砝码和木块相对桌面静止，由平衡条件可知 $f=\mu(M+m)g$，则有 $f=\mu Mg+\mu mg$，$f-m$ 图像中 $k=\mu g$。

（4）由 $k=\mu g$ 可得 $\mu=\frac{k}{g}=\frac{\frac{2.93-2.15}{0.25-0.05}}{9.80}\approx0.40$。
}

\item
%% number: 24
%% paperName: 2018年全国理综Ⅱ第 24 题 12 分
%% typeId: 6
%% type: 计算题
%% chapter: 运动和力
%% point: 动量守恒定律
%% method: 无
%% score: 12
%% degree: 650
%% duplicateId: 0
%% body:
汽车 A 在水平冰雪路面上行驶，驾驶员发现其正前方停有汽车 B，立即采取制动措施，但仍然撞上了汽车 B。两车碰撞时和两车都完全停止后的位置如图所示，碰撞后 B 车向前滑动了 $4.5\Um$，A 车向前滑动了 $2.0\Um$，已知 A 和 B 的质量分别为 $2.0\times10^{3}\Ukg$ 和 $1.5\times10^{3}\Ukg$，两车与该冰雪路面间的动摩擦因数均为 0.10，两车碰撞时间极短，在碰撞后车轮均没有滚动，重力加速度大小 $g=10\Umsq$。求：
\begin{enumerate}
	\item
	碰撞后的瞬间 B 车速度的大小；
	\item
	碰撞前的瞬间 A 车速度的大小。
\end{enumerate}

\onepicture[w=6.0cm, align=right]{figs/24.png}

%% answer:
\jdanswer{
\begin{enumerate}
	\item
	$v_{B}=3.0\Ums$

	\item
	$v_{A}=4.3\Ums$

\end{enumerate}
}
%% memo:
\memoanswer{
（1）设 B 车的质量为 $m_{B}$，碰后加速度大小为 $a_{B}$，由牛顿第二定律得 $\mu m_{B}g=m_{B}a_{B}$①，式中 $\mu$ 是汽车与路面间的动摩擦因数。设碰撞后瞬间 B 车速度的大小为 $v_{B}'$，碰撞后滑行的距离为 $s_{B}$，由运动学公式得 $v_{B}'^{2}=2a_{B}s_{B}$②，联立①②式并利用题给数据得 $v_{B}'=3.0\Ums$③。

（2）设 A 车的质量为 $m_{A}$，碰后加速度大小为 $a_{A}$，由牛顿第二定律得 $\mu m_{A}g=m_{A}a_{A}$④，设碰撞后瞬间 A 车速度的大小为 $v_{A}'$，碰撞后滑行的距离为 $s_{A}$，由运动学公式得 $v_{A}'^{2}=2a_{A}s_{A}$⑤，设碰撞前的瞬间 A 车速度的大小为 $v_{A}$。两车在碰撞过程中动量守恒，有 $m_{A}v_{A}=m_{A}v_{A}'+m_{B}v_{B}'$⑥，联立③④⑤⑥式并利用题给数据得 $v_{A}=4.3\Ums$⑦。
}

\item
%% number: 25
%% paperName: 2018年全国理综Ⅱ第 25 题 20 分
%% typeId: 6
%% type: 计算题
%% chapter: 磁场
%% point: 带电粒子在磁场中的运动
%% method: 无
%% score: 20
%% degree: 350
%% duplicateId: 0
%% body:
一足够长的条状区域内存在匀强电场和匀强磁场，其在 $xOy$ 平面内的截面如图所示：中间是磁场区域，其边界与 $y$ 轴垂直，宽度为 $l$，磁感应强度的大小为 $B$，方向垂直于 $xOy$ 平面；磁场的上、下两侧为电场区域，宽度均为 $l'$，电场强度的大小均为 $E$，方向均沿 $x$ 轴正方向；M、N 为条形区域边界上的两点，它们的连线与 $y$ 轴平行。一带正电的粒子以某一速度从 M 点沿 $y$ 轴正方向射入电场，经过一段时间后恰好以从 M 点入射的速度从 N 点沿 $y$ 轴正方向射出。不计重力。
\begin{enumerate}
	\item
	定性画出该粒子在电磁场中运动的轨迹；
	\item
	求该粒子从 M 点射入时速度的大小；
	\item
	若该粒子进入磁场时的速度方向恰好与 $x$ 轴正方向的夹角为 $\frac{\pi}{6}$，求该粒子的比荷及其从 M 点运动到 N 点的时间。
\end{enumerate}

\onepicture[w=5.5cm, align=right]{figs/25.png}

%% answer:
\jdanswer{
\begin{enumerate}
	\item
	轨迹图如图所示

	\item
	$v_{0}=\frac{2El'}{Bl}$

	\item
	$\frac{q}{m}=\frac{4\sqrt{3}El'}{B^{2}l^{2}}$，$t'=\frac{Bl}{E}\left(1+\frac{\sqrt{3}\pi l}{18l'}\right)$

\end{enumerate}
}
%% memo:
\memoanswer{
（1）粒子运动的轨迹如图(a)所示。（粒子在电场中的轨迹为抛物线，在磁场中为圆弧，上下对称）

\onepicture[w=5.0cm]{figs/25_an1.png}

（2）粒子从电场下边界入射后在电场中做类平抛运动。设粒子从 M 点射入时速度的大小为 $v_{0}$，在下侧电场中运动的时间为 $t$，加速度的大小为 $a$；粒子进入磁场的速度大小为 $v$，方向与电场方向的夹角为 $\theta$，如图所示，速度沿电场方向的分量为 $v_{1}$。由牛顿第二定律得 $qE=ma$①，式中 $q$ 和 $m$ 分别为粒子的电荷量和质量。由运动学公式得 $v_{1}=at$②，$l'=v_{0}t$③，$v_{1}=v\cos\theta$④，粒子在磁场中做匀速圆周运动，设其运动轨道半径为 $R$，由洛伦兹力公式和牛顿第二定律得 $qvB=\frac{mv^{2}}{R}$⑤，由几何关系得 $l=2R\cos\theta$⑥。

\onepicture[w=5.0cm]{figs/25_an2.png}

联立①②③④⑤⑥式得 $v_{0}=\frac{2El'}{Bl}$⑦。

（3）由运动学公式和题给数据得 $v_{1}=v_{0}\cot\frac{\pi}{6}$⑧，联立①②③⑦⑧式得 $\frac{q}{m}=\frac{4\sqrt{3}El'}{B^{2}l^{2}}$⑨，设粒子由 M 点运动到 N 点所用的时间为 $t'$，则 $t'=2t+\frac{2\left(\frac{\pi}{2}-\frac{\pi}{6}\right)}{2\pi}T$⑩，式中 $T$ 是粒子在磁场中做匀速圆周运动的周期，$T=\frac{2\pi m}{qB}$\ccnum{⑪}，由③⑦⑨⑩\ccnum{⑪}式得 $t'=\frac{Bl}{E}\left(1+\frac{\sqrt{3}\pi l}{18l'}\right)$\ccnum{⑫}。
}

\item
%% number: 33
%% paperName: 2018年全国理综Ⅱ第 33 题 10 分
%% typeId: 6
%% type: 计算题
%% chapter: 气体性质
%% point: 气体多过程
%% method: 无
%% score: 10
%% degree: 550
%% duplicateId: 0
%% body:
（1）对于实际的气体，下列说法正确的是\xzanswer{BDE}

\fivechoices[answer=BDE]
{气体的内能包括气体分子的重力势能}
{气体的内能包括分子之间相互作用的势能}
{气体的内能包括气体整体运动的动能}
{气体体积变化时，其内能可能不变}
{气体的内能包括气体分子热运动的动能}

（2）如图，一竖直放置的气缸上端开口，气缸壁内有卡口 a 和 b，a、b 间距为 $h$，a 距缸底的高度为 $H$；活塞只能在 a、b 间移动，其下方密封有一定质量的理想气体。已知活塞质量为 $m$，面积为 $S$，厚度可忽略；活塞和汽缸壁均绝热，不计它们之间的摩擦。开始时活塞处于静止状态，上、下方气体压强均为 $p_{0}$，温度均为 $T_{0}$。现用电热丝缓慢加热气缸中的气体，直至活塞刚好到达 b 处。求此时气缸内气体的温度以及在此过程中气体对外所做的功。重力加速度大小为 $g$。

\onepicture[h=4.0cm, align=right]{figs/33.png}

%% answer:
\jdanswer{
$T_{2}=\left(1+\frac{h}{H}\right)\left(1+\frac{mg}{p_{0}S}\right)T_{0}$，$W=(p_{0}S+mg)h$
}
%% memo:
\memoanswer{
（1）气体内能宏观上与气体的温度和体积有关，微观上与气体分子做热运动的平均动能和分子势能有关，对于实际气体，气体内能包括分子势能和分子动能，故 AC 错误、BE 正确；当气体体积变化时，若温度同时发生变化，气体内能可能不变，故 D 正确。

（2）开始时活塞位于 a 处，加热后，汽缸中的气体先经历等容过程，直至活塞开始运动。设此时汽缸中气体的温度为 $T_{1}$，压强为 $p_{1}$，由查理定律得 $\frac{p_{0}}{T_{0}}=\frac{p_{1}}{T_{1}}$①，由力的平衡条件得 $p_{1}S=p_{0}S+mg$②，联立①②式可得 $T_{1}=\left(1+\frac{mg}{p_{0}S}\right)T_{0}$③。此后，汽缸中的气体经历等压过程，直至活塞刚好到达 b 处，设此时汽缸中气体的温度为 $T_{2}$；活塞位于 a 处和 b 处时气体的体积分别为 $V_{1}$ 和 $V_{2}$，由盖-吕萨克定律得 $\frac{V_{1}}{T_{1}}=\frac{V_{2}}{T_{2}}$④，式中 $V_{1}=SH$⑤，$V_{2}=S(H+h)$⑥，联立③④⑤⑥式解得 $T_{2}=\left(1+\frac{h}{H}\right)\left(1+\frac{mg}{p_{0}S}\right)T_{0}$⑦。从开始加热到活塞到达 b 处的过程中，汽缸中的气体对外做的功为 $W=(p_{0}S+mg)h$⑧。
}

\item
%% number: 34
%% paperName: 2018年全国理综Ⅱ第 34 题 10 分
%% typeId: 6
%% type: 计算题
%% chapter: 光学
%% point: 光的折射
%% method: 无
%% score: 10
%% degree: 550
%% duplicateId: 0
%% body:
（1）声波在空气中的传播速度为 $340\Ums$，在钢铁中的传播速度为 $4900\Ums$。一平直桥由钢铁制成，某同学用锤子敲击一铁桥的一端而发出声音，分别经空气和桥传到另一端的时间之差为 $1.00\Us$。桥的长度为\tkanswer[1.6cm]{$365$} $\Um$，若该波在空气中的波长为 $\lambda$，则它在钢铁中波长为 $\lambda$ 的\tkanswer[1.8cm]{$\frac{245}{17}$}倍。

（2）如图，$\triangle ABC$ 是一直角三棱镜的横截面，$\angle A=90^{\circ}$，$\angle B=60^{\circ}$，一细光束从 BC 边的 D 点折射后，射到 AC 边的 E 点，发生全反射后经 AB 边的 F 点射出。EG 垂直于 AC 交 BC 于 G，D 恰好是 CG 的中点。不计多次反射。
\begin{enumerate}
	\item
	求出射光相对于 D 点的入射光的偏角；
	\item
	为实现上述光路，棱镜折射率的取值应在什么范围？
\end{enumerate}

\onepicture[w=6.0cm, align=right]{figs/34.png}

%% answer:
\jdanswer{
\begin{enumerate}
	\item
	$\delta=60^{\circ}$

	\item
	$\frac{2\sqrt{3}}{3}\leq n<2$

\end{enumerate}
}
%% memo:
\memoanswer{
（1）由于声音从桥的一端分别经空气和桥传播到另一端的时间差为 $1.00\Us$，则有 $\frac{L}{v_{\text{气}}}-\frac{L}{v_{\text{钢}}}=1.00\Us$，解得桥长为 $L=365\Um$；声音在不同介质中传播时，频率不变，故 $\frac{v_{\text{气}}}{\lambda_{\text{气}}}=\frac{v_{\text{钢}}}{\lambda_{\text{钢}}}$，可得 $\frac{\lambda_{\text{钢}}}{\lambda_{\text{气}}}=\frac{245}{17}$。

（2）(i) 光线在 BC 面上折射，由折射定律得 $\sin i_{1}=n\sin r_{1}$①，式中 $n$ 为棱镜的折射率，$i_{1}$ 和 $r_{1}$ 分别是该光线在 BC 面上的入射角和折射角。光线在 AC 面上发生全反射，由反射定律得 $i_{2}=r_{2}$②，式中 $i_{2}$ 和 $r_{2}$ 分别是该光线在 AC 面上的入射角和反射角。光线在 AB 面上发生折射，由折射定律得 $n\sin i_{3}=\sin r_{3}$③，式中 $i_{3}$ 和 $r_{3}$ 分别是该光线在 AB 面上的入射角和折射角。由几何关系得 $i_{2}=r_{2}=60^{\circ}$，$r_{1}=i_{3}=30^{\circ}$④，F 点的出射光相对于 D 点的入射光的偏角为 $\delta=(r_{1}-i_{1})+(180^{\circ}-i_{2}-r_{2})+(r_{3}-i_{3})$⑤，由①②③④⑤式得 $\delta=60^{\circ}$⑥。

（ii）光线在 AC 面上发生全反射，光线在 AB 面上不发生全反射，有 $n\sin i_{2}\geq n\sin C\geq n\sin i_{3}$⑦，式中 $C$ 是全反射临界角，满足 $n\sin C=1$⑧，由④⑦⑧式知，棱镜的折射率 $n$ 的取值范围应为 $\frac{2\sqrt{3}}{3}\leq n<2$⑨。

\onepicture[w=7.5cm]{figs/34_an.png}
}

\end{enumerate}

\end{document}
